%& -shell-escape
% !TeX root = thesis-text.tex
% !TeX encoding = UTF-8
% !TeX program = XeLaTeX

\documentclass{.local/lib/classes/ndvr-super-article-standalone}

\title[NDVR thesis-abstract]{%
    Проблема поиска нечетких дубликатов видео%
}

\begin{document}

    \pagebreak

    \section{Поиск шаблонов действий}

    За 20-30 лет человечество накопило большое 
    количество видео c камерах наружного наблюдения.
    Этот массив информации невозможно просмотреть целиком.
    Эффективных методов анализа видео, 
    подобных анализу текста, в настоящее время нет. 

    Для камер наружного наблюдения есть потребность искать 
    заданные шаблоны явлений в кадре. 
    Например, для систем безопасности актуальны шаблоны 
    «человек подходит к двери и открывает ее». 
    В реалистичной ситуации человек подходит 
    к двери с различных углов. 
    В кадре могут оказаться проезжающие мимо машины, 
    падающие листья, различные погодные явления.

    Практика показывает, что детекторы движения 
    будут ловить скорее листья и машины, 
    а не искомый шаблон поведения.

    Если в таком видео искать не движения, 
    а значимые события, то количество ложных срабатываний 
    можно свести к минимуму. Если в области обзора камеры 
    появился новый объект, которого раньше не было, 
    характеристики кадров видео изменятся. 
    Считаем, что произошло событие. 
    Если в области обзора оказались колышущиеся ветки 
    — то каждое движение веток тоже можно считать событием. 
    Но кадры с ветками до движения и во время движения 
    будут отличаться меньше. 
    Значимость события можно настроить 
    параметрами определения смены сцен. 

    Шаблоны «человек подходит к двери и открывает ее» 
    и «человек подходит к двери, стоит несколько секунд и уходит» 
    — разные шаблоны. Их нужно отличать. 
    Классические детекторы движения отличать такое не умеют.

    С видеокамерами связано еще одно приложение. 
    Последнее время стали широко применяться 
    бытовые БПЛА\index{БПЛА} («квадрокоптеры»\index{Квадрокоптер}). Обычно они оснащены камерами. 
    Такие малые летальные аппараты очень часто используются
    фотографами для создания снимков с удобных или неожиданных ракурсов. 
    Кроме того, «квадрокоптеры» используют путешественники 
    для проверки нового маршрута. 
    Похожие технологии начали применять 
    в геодезии и при составлении карт. 
    Для автоматической и полуавтоматической навигации 
    БПЛА тоже требуется поиск шаблонов. 
    Например, «пролетели над рекой, потом пролетели над дорогой». 

\end{document}